% !TeX program = pdflatex
\documentclass[10pt, aspectratio=169]{beamer}

\input{../../_en_preambulo_beamer}
% Comandos y macros estadísticas compartidas para las presentaciones Beamer en inglés (Metropolis)

% Conjuntos numéricos básicos
\newcommand{\R}{\mathbb{R}}
\newcommand{\N}{\mathbb{N}}
\newcommand{\Q}{\mathbb{Q}}
\newcommand{\Z}{\mathbb{Z}}
\newcommand{\set}[1]{\left\{ #1 \right\}}
\newcommand{\abs}[1]{\left|#1\right|}
\newcommand{\norm}[1]{\left\| #1 \right\|}

% Comandos estadísticos y probabilísticos del curso
\newcommand{\Var}{\operatorname{Var}}
\newcommand{\cov}{\operatorname{Cov}}
\newcommand{\corr}{\rho}
\newcommand{\comb}[2]{\begin{pmatrix} #1 \\ #2 \end{pmatrix}}
\newcommand{\s}{\sigma}
\newcommand{\particion}{\mathfrak{P}}
\newcommand{\card}[1]{\#\left( #1 \right)}
\newcommand{\seq}[3]{\set{#1}_{#2}^{#3}}
\newcommand{\E}{\mathbb{E}}
\newcommand{\Prob}{\mathbb{P}}

% Entornos pedagógicos y cajas en inglés académico natural (soportando tanto nombres en inglés como en español)
\newenvironment{discussion}{%
  \begin{alertblock}{Discussion Question}%
}{%
  \end{alertblock}%
}
\newenvironment{discusion}{%
  \begin{alertblock}{Discussion Question}%
}{%
  \end{alertblock}%
}

\newenvironment{reflection}{%
  \begin{alertblock}{Classroom Reflection}%
}{%
  \end{alertblock}%
}
\newenvironment{reflexion}{%
  \begin{alertblock}{Classroom Reflection}%
}{%
  \end{alertblock}%
}

\newenvironment{paradox}{%
  \begin{alertblock}{Exploratory Paradox}%
}{%
  \end{alertblock}%
}
\newenvironment{paradoja}{%
  \begin{alertblock}{Exploratory Paradox}%
}{%
  \end{alertblock}%
}

\newenvironment{motivation}{%
  \begin{exampleblock}{Motivation: Data Science in Practice}%
}{%
  \end{exampleblock}%
}
\newenvironment{motivacion}{%
  \begin{exampleblock}{Motivation: Data Science in Practice}%
}{%
  \end{exampleblock}%
}

\newenvironment{quickchallenge}{%
  \begin{block}{Quick Team Challenge}%
}{%
  \end{block}%
}
\newenvironment{retoclase}{%
  \begin{block}{Quick Team Challenge}%
}{%
  \end{block}%
}


\title[Random Sampling and CLT]{Section 02.06: Random Sampling and Central Limit Theorem}
\subtitle{From Brute-Force Census to Asymptotic Inference: SRS, LLN, and CLT}
\author[J. Castillo Colmenares]{Juliho Castillo Colmenares}
\institute{Tecnológico de Monterrey}
\date{\vspace{-2.0cm}}

\begin{document}

% Slide 01: Title
\begin{frame}[plain]
  \vspace{-0.5cm}
  \titlepage
\end{frame}

% Slide 02: Roadmap
\begin{frame}{Roadmap --- Chapter 02: Probability Theory}
  \begin{block}{Curricular Structure of Unit 1}
    \footnotesize
    \begin{enumerate}\itemsep=0.08cm
      \item \textcolor{gray}{Section 02.01: Introduction to Probability and Random Phenomena}
      \item \textcolor{gray}{Section 02.02: Set Theory and Kolmogorov's Axioms}
      \item \textcolor{gray}{Section 02.03: Fundamentals of Probability and Counting Techniques}
      \item \textcolor{gray}{Section 02.04: Conditional Probability and Statistical Independence}
      \item \textcolor{gray}{Section 02.05: Bayes' Theorem and Conditional Inversion}
      \item \textbf{\textcolor{TecRojo}{Section 02.06: Random Sampling and Central Limit Theorem (CLT)}} $\leftarrow$ \emph{Current focus}
    \end{enumerate}
  \end{block}
  \vspace{-0.05cm}
  \pause
  \begin{alertblock}{Didactic Goal of Section 02.06}
    \footnotesize
    To master the rigorous foundations of simple random sampling (SRS), quantify the finite population correction factor ($\text{FPC}$), and command the asymptotic laws of sample statistics: the Law of Large Numbers (LLN) and the Central Limit Theorem (CLT).
  \end{alertblock}
\end{frame}

% Slide 03: Motivation — From Brute-Force Census to Random Sampling
\begin{frame}{Motivation --- From Brute-Force Census to Random Sampling}
  \begin{columns}[T]
    \begin{column}{0.48\textwidth}
      \begin{block}{The Brute-Force Dilemma (Census)}
        \small
        \begin{itemize}\itemsep=0.1cm
          \item In engineering, economics, and life sciences, determining exact parameters across massive populations by measuring every unit is infeasible:
          \begin{itemize}
            \item Prohibitive logistical and computational costs.
            \item Destructive quality testing (e.g., tensile strength of steel beams).
            \item Conceptually infinite target populations.
          \end{itemize}
        \end{itemize}
      \end{block}
    \end{column}
    \begin{column}{0.48\textwidth}
      \begin{alertblock}{The Statistical Solution: SRS Inference}
        \small
        \begin{itemize}\itemsep=0.1cm
          \item We extract a representative subset of size $n \ll N$ under rigorous conditions of randomness.
          \item We compute a \textbf{sample statistic} ($\bar{X}, S^2$) to approximate and bound the \textbf{population parameter} ($\mu, \sigma^2$) with known probability.
          \item We transform qualitative uncertainty into quantifiable, bounded statistical risk.
        \end{itemize}
      \end{alertblock}
    \end{column}
  \end{columns}
\end{frame}

% Slide 04: Simple Random Sampling (SRS) and Statistical Independence
\begin{frame}{Simple Random Sampling (SRS) and Statistical Independence}
  \fontsize{7.5pt}{9pt}\selectfont
  \begin{block}{Formal Definition of Simple Random Sampling (SRS)}
    \fontsize{7.5pt}{9pt}\selectfont
    A collection of random variables $X_1, X_2, \ldots, X_n$ constitutes a \textbf{simple random sample of size $n$} drawn from the population $f_X(x)$ if they simultaneously fulfill two axiomatic conditions:
    \begin{enumerate}\itemsep=0.1cm
      \item \textbf{Identical Distribution:} Each individual variable $X_i$ shares the exact same marginal probability distribution, density, and theoretical moments as $X$: $f_{X_i}(x) = f_X(x)$, with $E(X_i) = \mu$ and $\text{Var}(X_i) = \sigma^2$.
      \item \textbf{Mutual Independence:} The observation of any subset of variables does not alter the conditional distribution of the remaining variables. The joint probability density factors perfectly:
      $$f_{X_1, X_2, \ldots, X_n}(x_1, x_2, \ldots, x_n) = \prod_{i=1}^n f_X(x_i)$$
    \end{enumerate}
  \end{block}
  \vspace{-0.15cm}
  \pause
  \begin{alertblock}{Statistics vs. Parameters}
    \fontsize{7.5pt}{9pt}\selectfont
    A \textbf{parameter} is a fixed but unknown real constant of the population ($\mu, \sigma^2$). A \textbf{statistic} is a measurable function of the random sample $T(X_1, \ldots, X_n)$ involving no unknown constants ($\bar{X}, S^2$). Being a function of random variables, every statistic is itself a random variable with its own probability distribution, known as its \emph{sampling distribution}.
  \end{alertblock}
\end{frame}

% Slide 05: Sampling With vs Without Replacement (FPC Factor)
\begin{frame}{Sampling With vs Without Replacement (FPC Factor)}
  \fontsize{7.5pt}{9pt}\selectfont
  \begin{block}{Expected Value and Variance in Infinite Populations (or With Replacement)}
    \fontsize{7.5pt}{9pt}\selectfont
    If sampling is performed with replacement or from an infinite population ($N \to \infty$), the variables $X_i$ are strictly i.i.d. By linearity of expectation and additivity of variance for independent variables:
    $$E(\bar{X}) = E\left[\frac{1}{n}\sum_{i=1}^n X_i\right] = \frac{1}{n}(n\mu) = \mu, \qquad \text{Var}(\bar{X}) = \text{Var}\left[\frac{1}{n}\sum_{i=1}^n X_i\right] = \frac{1}{n^2}(n\sigma^2) = \frac{\sigma^2}{n}$$
  \end{block}
  \vspace{-0.15cm}
  \pause
  \begin{alertblock}{The Finite Population Correction (FPC) Factor}
    \fontsize{7.5pt}{9pt}\selectfont
    When sampling without replacement from a finite population of size $N$, the variables $X_i$ exhibit a slight negative correlation induced by census depletion ($\text{Cov}(X_i, X_j) = -\frac{\sigma^2}{N-1}$). This dependency strictly reduces the sample mean variance via the \textbf{Finite Population Correction (FPC) factor}:
    $$\text{Var}_{W/O}(\bar{X}) = \frac{\sigma^2}{n} \cdot \left(\frac{N-n}{N-1}\right) \implies \sigma_{\bar{X}} = \frac{\sigma}{\sqrt{n}} \sqrt{\frac{N-n}{N-1}}$$
    If the sampling fraction $n/N \le 0.05$ ($5\%$), the factor $(N-n)/(N-1) \approx 1$ and the FPC is safely omitted in practice.
  \end{alertblock}
\end{frame}

% Slide 06: The Law of Large Numbers (LLN)
\begin{frame}{The Law of Large Numbers (LLN)}
  \fontsize{7.5pt}{9pt}\selectfont
  \begin{block}{Weak Law of Large Numbers (WLLN — Khinchin)}
    \fontsize{7.5pt}{9pt}\selectfont
    Let $X_1, X_2, \ldots, X_n$ be an i.i.d. random sample drawn from a population with finite theoretical mean $\mu < \infty$. For any arbitrarily small margin of tolerance $\epsilon > 0$, the probability that the sample mean $\bar{X}_n$ deviates from the population parameter $\mu$ by $\epsilon$ or more tends rigorously to zero as the sample size $n$ increases:
    $$\lim_{n \to \infty} P(|\bar{X}_n - \mu| \ge \epsilon) = 0 \quad \iff \quad \bar{X}_n \xrightarrow{p} \mu \quad \text{(Convergence in Probability)}$$
  \end{block}
  \vspace{-0.15cm}
  \pause
  \begin{alertblock}{Strong Law of Large Numbers (SLLN — Kolmogorov)}
    \fontsize{7.5pt}{9pt}\selectfont
    Under the exact same assumption of a finite first-order moment, the sequence of sample means converges to the true population parameter $\mu$ with probability exactly equal to one:
    $$P\left(\lim_{n \to \infty} \bar{X}_n = \mu\right) = 1 \quad \iff \quad \bar{X}_n \xrightarrow{a.s.} \mu \quad \text{(Almost Sure Convergence)}$$
    \textbf{Operational Meaning:} Empirical uncertainty regarding the center of gravity of a random phenomenon vanishes asymptotically when accumulating sufficient independent observations.
  \end{alertblock}
\end{frame}

% Slide 07: The Central Limit Theorem (CLT)
\begin{frame}{The Central Limit Theorem (CLT)}
  \fontsize{7.5pt}{9pt}\selectfont
  \begin{block}{Classic Lindeberg-Lévy Statement}
    \fontsize{7.5pt}{9pt}\selectfont
    Let $X_1, X_2, \ldots, X_n$ be a simple random sample (i.i.d.) drawn from a population with theoretical mean $\mu$ and strictly positive, finite variance $0 < \sigma^2 < \infty$. As the sample size $n$ grows indefinitely, the standardized random variable $Z_n$ associated with the sample mean converges in distribution to the standard normal distribution $N(0,1)$, \textbf{regardless of the shape or skewness of the underlying population distribution}:
    $$\lim_{n \to \infty} P\left(Z_n \le z\right) = \lim_{n \to \infty} P\left(\frac{\bar{X}_n - \mu}{\sigma / \sqrt{n}} \le z\right) = \Phi(z) = \frac{1}{\sqrt{2\pi}} \int_{-\infty}^{z} e^{-t^2/2} \, dt$$
  \end{block}
  \vspace{-0.15cm}
  \pause
  \begin{alertblock}{The Universal Scope of the CLT}
    \fontsize{7.5pt}{9pt}\selectfont
    \begin{itemize}\itemsep=0.08cm
      \item \textbf{Asymptotic Approximation of the Mean:} For sufficiently large $n$ ($n \ge 30$ as a standard empirical rule of thumb), the sample mean can be treated as normal: $\bar{X}_n \dot{\sim} N\left(\mu, \frac{\sigma^2}{n}\right)$.
      \item \textbf{Approximation of the Sample Sum:} Analogously, the total sum $S_n = \sum_{i=1}^n X_i$ converges to a normal distribution: $S_n \dot{\sim} N(n\mu, n\sigma^2)$.
    \end{itemize}
  \end{alertblock}
\end{frame}

% Slide 08: Asymptotic Anatomy and Berry-Esseen Bound
\begin{frame}{Asymptotic Anatomy and Rate of Convergence}
  \begin{columns}[T]
    \begin{column}{0.48\textwidth}
      \begin{block}{The Empirical $n \ge 30$ Rule}
        \small
        \begin{itemize}\itemsep=0.1cm
          \item If the original population is normal, $\bar{X}_n \sim N(\mu, \sigma^2/n)$ exactly for any $n \ge 1$.
          \item If the population is uniform, symmetric, or continuous with finite variance, $n \ge 15$ usually suffices for excellent normality.
          \item If the population is highly skewed (exponential, Pareto), $n \ge 30$ to $n \ge 100$ is required to stabilize tail probabilities.
        \end{itemize}
      \end{block}
    \end{column}
    \begin{column}{0.48\textwidth}
      \begin{alertblock}{Universal Berry-Esseen Bound}
        \small
        \begin{itemize}\itemsep=0.1cm
          \item To mathematically bound the maximum discrepancy between the empirical CDF $F_n(z)$ and $\Phi(z)$, the absolute third central moment $\rho_3 = E[|X - \mu|^3] < \infty$ is required:
          $$\sup_{z \in \mathbb{R}} |F_n(z) - \Phi(z)| \le \frac{C \cdot \rho_3}{\sigma^3 \sqrt{n}}$$
          \item The approximation error decays at rate $O(1/\sqrt{n})$, where $C < 0.4784$.
        \end{itemize}
      \end{alertblock}
    \end{column}
  \end{columns}
\end{frame}

% Slide 09A: Python Lab — Sampling With vs Without Replacement
\begin{frame}[fragile]{Python Lab --- Sampling With vs Without Replacement (1/4)}
  \vspace{-0.1cm}
  \begin{block}{Source Code: \texttt{compare\_sampling\_schemes} (\texttt{02.06\_random\_sampling.py})}
    \lstinputlisting[language=Python, basicstyle=\fontsize{4.6pt}{5.4pt}\selectfont\ttfamily, firstline=10, lastline=28]{../../code/02_teoria_probabilidad/02.06_random_sampling.py}
  \end{block}
\end{frame}

% Slide 09B: Python Lab — Law of Large Numbers
\begin{frame}[fragile]{Python Lab --- Law of Large Numbers (2/4)}
  \vspace{-0.1cm}
  \begin{block}{Source Code: \texttt{verify\_law\_of\_large\_numbers} (\texttt{02.06\_random\_sampling.py})}
    \lstinputlisting[language=Python, basicstyle=\fontsize{5.5pt}{6.5pt}\selectfont\ttfamily, firstline=32, lastline=46]{../../code/02_teoria_probabilidad/02.06_random_sampling.py}
  \end{block}
\end{frame}

% Slide 09C: Python Lab — Central Limit Theorem
\begin{frame}[fragile]{Python Lab --- Central Limit Theorem (3/4)}
  \vspace{-0.1cm}
  \begin{block}{Source Code: \texttt{verify\_central\_limit\_theorem} (\texttt{02.06\_random\_sampling.py})}
    \lstinputlisting[language=Python, basicstyle=\fontsize{4.6pt}{5.4pt}\selectfont\ttfamily, firstline=50, lastline=72]{../../code/02_teoria_probabilidad/02.06_random_sampling.py}
  \end{block}
\end{frame}

% Slide 09D: Python Lab — Verified Console Output
\begin{frame}[fragile]{Python Lab --- Verified Console Output (4/4)}
  \vspace{-0.1cm}
  \begin{columns}[T]
    \begin{column}{0.54\textwidth}
      \begin{block}{1. SRS With vs Without Replacement \& 2. LLN}
        \fontsize{4.4pt}{5.2pt}\selectfont
\begin{verbatim}
--- 1. Sampling With vs Without Replacement ---
Population N = 500, Sample n = 50, True Sigma = 19.6054
Empirical Var(X_bar) With Replacement    : 7.6989
Empirical Var(X_bar) Without Replacement : 7.0123
Theoretical FPC Factor [(N-n)/(N-1)]     : 0.9018
Empirical Variance Ratio (Without / With): 0.9108

--- 2. Law of Large Numbers (LLN) Convergence ---
Sample Size (n) | Empirical Mean | Absolute Error
-------------------------------------------------
             10 |         3.8000 |         0.3000
             50 |         3.5400 |         0.0400
            200 |         3.4400 |         0.0600
           1000 |         3.4990 |         0.0010
           5000 |         3.5200 |         0.0200
\end{verbatim}
      \end{block}
    \end{column}
    \begin{column}{0.44\textwidth}
      \begin{block}{3. CLT Standardization}
        \fontsize{4.4pt}{5.2pt}\selectfont
\begin{verbatim}
--- 3. Central Limit Theorem (CLT) ---
Population: Exponential(mu=2, sigma=2), n=36
Empirical P(|Z| <= 1.00) : 0.6821
Theoretical N(0,1)       : 0.6827

Empirical P(|Z| <= 1.96) : 0.9510
Theoretical N(0,1)       : 0.9500

Empirical P(|Z| <= 2.00) : 0.9556
Theoretical N(0,1)       : 0.9545
\end{verbatim}
      \end{block}
    \end{column}
  \end{columns}
\end{frame}

% Slide 14: Synthesis and Conclusion
\begin{frame}{Synthesis and Conclusion --- Random Sampling and CLT}
  \fontsize{7.5pt}{9pt}\selectfont
  \begin{block}{Conceptual Pillars of Section 02.06}
    \fontsize{7.5pt}{9pt}\selectfont
    \begin{itemize}\itemsep=0.12cm
      \item \textbf{Simple Random Sampling (SRS):} An i.i.d. random sample $X_1, \ldots, X_n$ enables the construction of unbiased sample statistics to bound uncertainty concerning population parameters.
      \item \textbf{FPC Factor in Finite Populations:} Sampling without replacement across a census fraction strictly reduces estimator variance: $\text{FPC} = \frac{N-n}{N-1}$.
      \item \textbf{Asymptotic Stability (LLN):} The Law of Large Numbers guarantees that the sample mean $\bar{X}_n$ converges in probability (WLLN) and almost surely (SLLN) toward $\mu$.
      \item \textbf{Universal Normality (CLT):} Provided $\sigma^2 < \infty$, standardized sample means $Z_n$ converge to $N(0,1)$ for large sample sizes ($n \ge 30$), irrespective of initial population skewness.
    \end{itemize}
  \end{block}
\end{frame}

% Slide 15: Curricular Transition and Next Topic
\begin{frame}{Curricular Transition to the Next Section}
  \fontsize{7.5pt}{9pt}\selectfont
  \begin{block}{Looking Backward and Looking Forward}
    \fontsize{7.5pt}{9pt}\selectfont
    With the formalization of Random Sampling and Limit Theorems, we conclude the foundational architecture of \textbf{Classical Probability Theory (Chapter 02)}.
    \begin{itemize}\itemsep=0.1cm
      \item Empirical sampling provides us with sample evidence (likelihood) to update prior knowledge regarding a random phenomenon.
      \item To formally merge sample data with prior probability distributions, we require the powerful analytical engine of \textbf{conditional probability inversion}.
    \end{itemize}
  \end{block}
  \vspace{-0.15cm}
  \pause
  \begin{alertblock}{Next Stop in Curricular Remediation: Section 02.05}
    \fontsize{7.5pt}{9pt}\selectfont
    Moving forward in our systematic remediation plan, we will transition to \textbf{Section 02.05: Bayes' Theorem and Conditional Inversion}, where we explore how the Law of Total Probability combines with the product rule to compute posterior probabilities $P(A_i \mid B)$ in technical diagnosis and engineering systems.
  \end{alertblock}
\end{frame}

\end{document}
